% linear_probing.tex
\begin{figure}[H]
	\centering
	\begin{tikzpicture}[
		slot/.style={draw, minimum width=3.5cm, minimum height=0.8cm, align=center},
		occupied/.style={slot, fill=gray!30},
		empty/.style={slot, fill=white},
		>=Stealth,
		start chain=1 going right,
		node distance=-\pgflinewidth
		]
		
		% Array structure
		\node [on chain=1] {\dots};
		\node [on chain=1, occupied] (slot0) {Occupied};
		\node [on chain=1, occupied] (slot1) {Occupied};
		\node [on chain=1, empty]    (slot2) {Empty};
		\node [on chain=1] {\dots};
		
		% Indices below the array
		\node[below=0.3cm of slot0] {$h(k) \bmod M$};
		\node[below=0.3cm of slot1] {$h(k)+1 \bmod M$};
		\node[below=0.3cm of slot2] {$h(k)+2 \bmod M$};
		
		% Key input
		\node[above=2.5cm of slot0, draw, rounded corners, inner sep=0.2cm] (key) {Insert Key $k$};
		
		% Probing Arrows
		\draw[->, thick] (key.south) -- (slot0.north)
		node[midway, left=0.1cm] {$i=0$};
		
		\draw[->, thick] (slot0.north) .. controls +(up:1.5cm) and +(up:1.5cm) .. (slot1.north)
		node[midway, above] {$i=1$ (Collision)};
		
		\draw[->, thick] (slot1.north) .. controls +(up:1.5cm) and +(up:1.5cm) .. (slot2.north)
		node[midway, above] {$i=2$ (Success)};
		
	\end{tikzpicture}
	\caption{Pictorial representation of Linear Probing in Hash Tables}
	\label{fig:linear_probing}
\end{figure}